\documentclass{article}
% ML4PS 2026 (NeurIPS workshop). 4 pages excluding references.
% Deadline 12 September 2026, 23:59 AoE.
% Build: the neurips_2026.sty file is not vendored here -- fetch it from the
% workshop site and drop it beside this file, or build on Overleaf.
\usepackage[final]{neurips_2026}

\usepackage[utf8]{inputenc}
\usepackage[T1]{fontenc}
\usepackage{hyperref}
\usepackage{url}
\usepackage{booktabs}
\usepackage{amsmath}
\usepackage{amssymb}
\usepackage{graphicx}
\usepackage{xcolor}
\usepackage{siunitx}

% ============================================================================
% SINGLE SOURCE OF TRUTH FOR EVERY NUMBER IN THE PAPER.
%
% Nothing in main.tex may hardcode a result. Each macro below carries the
% run_id it came from, so any figure in the text can be traced to a row of
% experiments/RUNS.csv without opening another file.
%
% \pending{} renders in red. A compiled draft with red in it is not submittable
% -- that is the point. Every \pending must be resolved before 12 Sep.
% ============================================================================
\newcommand{\pending}[1]{\textcolor{red}{\textbf{[#1]}}}

% ---- dataset and model (docs/GROUND_TRUTH.md) ----
\newcommand{\Nnative}{188}
\newcommand{\Nresonant}{161}
\newcommand{\Nqcd}{27}
\newcommand{\Ntrainjets}{110{,}087{,}593}
\newcommand{\Ntestjets}{27{,}448{,}839}
\newcommand{\Nparams}{2{,}304{,}688}
\newcommand{\Nexamples}{8.192\times10^{8}}
\newcommand{\LR}{5\times10^{-4}}
\newcommand{\Nepochs}{80}

% ---- G1 learning-rate sweep (run_id g1-*) ----
\newcommand{\tiemargin}{0.0158}
\newcommand{\LsixtwoLRa}{0.52183}   % 2.5e-4
\newcommand{\LsixtwoLRb}{0.55232}   % 5e-4
\newcommand{\LsixtwoLRc}{0.55548}   % 1e-3
\newcommand{\RsixteenLRa}{0.77917}
\newcommand{\RsixteenLRb}{0.78277}
\newcommand{\RsixteenLRc}{0.75986}
\newcommand{\lrgap}{0.00316}        % L162: 1e-3 minus 5e-4

% ---- pretraining loss, training logs 2026-09-03 ----
% BOTH AT EPOCH 79 OF A COMPLETED RUN. The first draft of this table took
% R16_Q1 from mtx-r16q1-s1, which is at epoch 29 and still running -- an
% unfair comparison against L162 at epoch 79. Sources are now
% mtx-l162-s1b (80/80) and mtx-r16q1-s5 (80/80).
\newcommand{\LossLsixtwo}{1.599}    % mtx-l162-s1b, epoch 79
\newcommand{\LossRsixteen}{0.633}   % mtx-r16q1-s5, epoch 79
\newcommand{\ChanceLsixtwo}{5.088}  % ln 162
\newcommand{\ChanceRsixteen}{2.833} % ln 17
\newcommand{\NormLsixtwo}{0.314}
\newcommand{\NormRsixteen}{0.223}

% ---- frozen probe. v1 = 20% budget, CONFOUNDED. v2 = full budget. ----
% v1 (run_id probe-bvc-v1) -- reported only as the preliminary in the appendix,
% never as the headline: its arms ran at different learning rates.
\newcommand{\vaLsixtwoLin}{0.99689}
\newcommand{\vaLsixtwoMlp}{0.99698}
\newcommand{\vaRsixteenLin}{0.98813}
\newcommand{\vaRsixteenMlp}{0.98881}
\newcommand{\vaDlogLin}{-1.3405}
\newcommand{\vaDlogLinCI}{[-1.5585,\,-1.1438]}
\newcommand{\vaDlogMlp}{-1.3111}
\newcommand{\vaDlogMlpCI}{[-1.5192,\,-1.1247]}
\newcommand{\vaNjets}{60{,}565}

% v2 (run_id probe-bvc-v2) -- THE HEADLINE. Fill from
% /data/results/eval/probe_bvc_v2/probe_results.json when the job returns.
\newcommand{\vbLsixtwoLin}{\pending{v2 L162 linear}}
\newcommand{\vbLsixtwoMlp}{\pending{v2 L162 mlp}}
\newcommand{\vbRsixteenLin}{\pending{v2 R16 linear}}
\newcommand{\vbRsixteenMlp}{\pending{v2 R16 mlp}}
\newcommand{\vbDlogLin}{\pending{v2 dlog linear}}
\newcommand{\vbDlogLinCI}{\pending{v2 CI}}
\newcommand{\vbDlogMlp}{\pending{v2 dlog mlp}}
\newcommand{\vbDlogMlpCI}{\pending{v2 CI}}
\newcommand{\vbNjets}{\pending{v2 n}}
\newcommand{\vbSeedSpread}{\pending{R16 seed spread}}

% retained_topology -- the CONTROL leg. Both arms keep this axis, so the
% L162-minus-R16 gap here should be ~0 while the gap on bvc_resonant is large.
% That difference-in-differences is what separates "this axis was erased" from
% "this arm is worse at everything".
\newcommand{\rtLsixtwoLin}{\pending{retained L162 linear}}
\newcommand{\rtLsixtwoMlp}{\pending{retained L162 mlp}}
\newcommand{\rtRsixteenLin}{\pending{retained R16 linear}}
\newcommand{\rtRsixteenMlp}{\pending{retained R16 mlp}}
\newcommand{\rtDlogLin}{\pending{retained dlog}}
\newcommand{\rtDlogLinCI}{\pending{retained CI}}

% bvc_qcd -- the NULL-vs-NULL leg. BOTH arms collapse it, so neither should
% win. A gap here would indict the method, not the vocabulary.
\newcommand{\qqLsixtwoLin}{\pending{qcd L162 linear}}
\newcommand{\qqRsixteenLin}{\pending{qcd R16 linear}}
\newcommand{\qqDlogLin}{\pending{qcd dlog}}
\newcommand{\qqDlogLinCI}{\pending{qcd CI}}

% ---- pipeline validation ----
\newcommand{\ePartOurs}{0.9878}
\newcommand{\ePartPub}{0.9877}
\newcommand{\sigmapretrain}{0.003704}


\title{What must a pretraining vocabulary preserve?\\
       A controlled label-granularity ladder for jet foundation models}

\author{%
  Raunav Mendiratta\thanks{Correspondence: \texttt{raunav.mendiratta@gmail.com}} \\
  % TODO: full author list and affiliations before submission
}

\begin{document}
\maketitle

\begin{abstract}
Jet foundation models are pretrained on simulated collision data with labels
whose granularity is chosen by convention rather than measured. Two axes of
pretraining design have recently been ablated --- the training objective and the
composition of the pretraining data --- but not the label space. We pretrain Particle Transformers on JetClass-II
under a nested contraction of its \Nnative{} native labels, holding the data
stream's sampling distribution byte-identical across arms so that the
vocabulary is the only deliberate variable.
Probing the frozen representation on three axes that differ in how the
vocabulary touches them, we find a clean ordering. Coarsening degrades the
representation generally, by a factor $\facTop\times$ in $1-\mathrm{AUC}$ even
on an axis both arms resolve. A distinction the coarse arm merges degrades a
further $\facRes\times$. And an axis that \emph{neither} arm labels degrades by
$\facQcd\times$: the fine arm decodes QCD flavour having never received a QCD
flavour label, and a single fitted direction serves both flavour axes, so this is
the same erased feature read on jets whose labels never carried it. What a
vocabulary erases is not confined to the jets whose labels it merged. Whether the
coarse arm loses it because those labels were merged or because $K=17$ compresses
every axis interior to a group is what our pre-registered random-partition
control separates. Pretraining loss, separately, is defined relative to a label
space: our arms' losses are not the same quantity, and which arm learned more is
not decidable from them.
\end{abstract}

\section{Introduction}

Foundation models for jet physics are now trained at a scale where the
pretraining recipe, not the architecture, is the object of study. Two design
axes have recently been ablated under controlled conditions. \citet{objectives}
compare pretraining \emph{objectives} and find pure classifier pretraining
optimal when downstream labels and capacity are plentiful; \citet{datacomp} vary
the \emph{composition} of the pretraining data and show that scaling behaviour
can be engineered by making the corpus more diverse and better aligned with the
downstream task.

A third axis has not been ablated: the \emph{label space}. If supervised
classification is the best pretraining objective, then the vocabulary of that
classification is the design choice that matters most, and it is currently set
by convention. Production taggers have grown from tens to hundreds of classes
without a controlled measurement of what that buys, and the vision literature
disagrees with itself about the answer: \citet{huh} found \emph{coarsening}
ImageNet helped transfer to few-class targets, \citet{hong} found finer labels
monotonically better on ImageNet21k but U-shaped on iNaturalist, and
\citet{mehta} report saturation at roughly eight categories.

We contribute a controlled instrument for that question in jet physics: our
ladder draws every arm from an identical training distribution, so the
vocabulary is the only deliberate variable --- a control none of those studies
could apply, since coarsening an image taxonomy changes which images are
sampled.

\section{A controlled granularity ladder}

\paragraph{Data and vocabulary.} JetClass-II labels each jet with one of
\Nnative{} native classes: \Nresonant{} resonant decay signatures and \Nqcd{}
QCD categories. We define a nested contraction tree over these labels and study
two rungs of it here: \textbf{L162}, which merges only the QCD block
($K=162$), and \textbf{R16\_Q1}, which merges the resonant signatures into
16 topological groups plus one QCD node ($K=17$). The map is materialised, not
regenerated from a seed, so every arm sees an identical label array.

\paragraph{What is held fixed.} All arms use the same Particle Transformer
trunk \citep{part}; only the output layer's width follows $K$, giving
$\NparamsLsixtwo$ trainable parameters for L162 against $\NparamsRsixteen$ for
R16\_Q1. Every arm sees the same $\Nexamples$ examples over \Nepochs{} epochs at the
same learning rate $\LR$, and selection retains \Ntrainjets{} training and
\Ntestjets{} test jets. Crucially, the flat $(p_T, m_\mathrm{SD})$ reweighting
histograms that determine \emph{which jets are drawn} are verified byte-identical
across arms, so the sampling distribution, selection and file list are common to
both. The seed is not: L162 runs at seed 1 and R16\_Q1 at seeds 2--4, so draw
order, trunk initialisation and dropout masks differ between arms as well as
between seeds. This is one seed against three rather than a matched pair, and the
seed-level tests of Section~3.1 are what carry it.

\paragraph{One learning rate is defensible.} A separate rate per arm would
confound granularity with tuning. Sweeping three rates per arm at 20\% budget,
the seed spread at fixed rate is \tiemargin{} in validation accuracy, which sets
the resolution. At $\LR$ R16\_Q1 attains its outright optimum (\RsixteenLRb) and
L162 sits \lrgap{} below its own (\LsixtwoLRb{} against \LsixtwoLRc{}) --- one
fifth of the noise floor, i.e.\ a tie.

\paragraph{Endpoint.} Arms with different $K$ emit incommensurable logits, so
we compare the frozen 128-dimensional representation the classifier head reads,
probed on tasks defined over the \emph{native} labels and therefore identical for
every arm. We fit both a linear and a nonlinear (MLP) probe, and never report one
without the other.

\section{Results}

\subsection{Coarsening damages an axis it never labelled}

A single contrast on a collapsed axis cannot support the claim, because a
coarser vocabulary might simply yield a weaker representation of everything. We
therefore probe three axes that differ in how the vocabulary touches them, and
read the arms against each other \emph{within} an axis.

\textbf{Flavour, directly supervised.} $X\!\to\!b\bar b$ against
$X\!\to\!c\bar c$, both two-prong hadronic, so topology is held fixed. L162
separates these two native labels; R16\_Q1 merges them.
\textbf{Topology, kept by both.} $X\!\to\!b\bar b$ against
$X\!\to\!YY\!\to\!b\bar bb\bar b$: flavour held fixed (both endpoints all-$b$),
prong count varied. Both arms separate these.
\textbf{Flavour, supervised in neither arm.} QCD $b\bar b$ against QCD
$c\bar c$. Both arms merge all 27 QCD classes into a single node, so neither
receives any QCD flavour label. We ran three R16\_Q1 pretraining seeds against
one L162 seed.

\begin{table}[h]
\centering
\caption{$\Delta\log(1-\mathrm{AUC})$ of each R16\_Q1 seed against L162, frozen
linear probe, 95\% paired-bootstrap CI. Arms are scored on identical jets in
identical order. Negative means L162 is better. The three columns are ordered
by how directly the vocabulary touches the axis. Intervals resample the jets of
the held-out probe split ($n=\nTestRes/\nTestQcd/\nTestTop$, 2000 resamples)
conditional on one checkpoint per arm and one probe fit; they carry no
pretraining-seed and no probe-training component. The mean and factor rows are
point estimates; their seed-level uncertainty is given in the text.}
\label{tab:probe}
\begin{tabular}{lccc}
\toprule
 & flavour & flavour & topology \\
 & (L162 labels it) & (neither labels it) & (both label it) \\
\midrule
seed 2 & \vbDlogA\ \vbDlogACI & \qqDlogA\ \qqDlogACI & \rtDlogA\ \rtDlogACI \\
seed 3 & \vbDlogB\ \vbDlogBCI & \qqDlogB\ \qqDlogBCI & \rtDlogB\ \rtDlogBCI \\
seed 4 & \vbDlogC\ \vbDlogCCI & \qqDlogC\ \qqDlogCCI & \rtDlogC\ \rtDlogCCI \\
\midrule
mean & \meanRes & \meanQcd & \meanTop \\
factor in $1-\mathrm{AUC}$ & $\facRes\times$ & $\facQcd\times$ & $\facTop\times$ \\
\bottomrule
\end{tabular}
\end{table}

\paragraph{Two effects, cleanly ordered.} The coarse arm is worse even on the
axis it \emph{keeps} --- a factor $\facTop\times$. Coarsening degrades the
representation generally, and a study without that column would have charged
this cost to the erased distinction. Since that axis is one the coarse arm still
labels, $\facTop\times$ is a \emph{floor}, and the excesses below are upper
bounds on the vocabulary-specific component. The axis R16\_Q1 collapses loses a
further
$\didRes$ in $\log(1-\mathrm{AUC})$ on top of it. For every seed the three
intervals are mutually disjoint, so the ordering holds within each checkpoint
without a bootstrap of the differences. Subtracting this way assumes general
degradation is additive in $\log(1-\mathrm{AUC})$ across axes of unequal
difficulty; calibrating a uniform $d'$ loss on the retained axis instead gives
$\didResDprime$ here and $\didQcdDprime$ on the unlabelled axis. Both excesses
stay positive under either model and move in opposite directions between them,
so the ordering is robust and the magnitudes are model-dependent.

\paragraph{The damage reaches an axis the vocabulary never labelled.}
\emph{Neither} arm receives a QCD flavour label: both merge all 27 QCD classes
into one node, so the loss is flavour-invariant on QCD jets in both arms. We
designed this column as a null-versus-null control and expected parity; that
prior was wrong, and in hindsight it cut against the routine transfer of flavour
taggers across jet populations. L162 is better by $\facQcd\times$, an excess of
$\didQcd$ over the retained axis.

The two flavour axes are one feature, not two. A single direction fitted on
both jointly reaches essentially in-domain AUC on each ($\poolLsixtwoRes$ and
$\poolLsixtwoQcd$ for L162, against $\inLsixtwoRes$ and $\inLsixtwoQcd$
in-domain; likewise for every R16\_Q1 seed), and the direction fitted on QCD jets
alone --- jets that carried no flavour label in either arm --- separates resonant
$b$ from $c$ at AUC $\xferQcdToRes$ in L162. So the QCD column is not a third
effect of a different kind but the erased feature read on jets whose labels never
carried it, which is what makes the merged-axis result more than memorisation of
a split between two training labels. One
explanation is that flavour discrimination learned from the \emph{resonant}
labels transfers to QCD jets, and that collapsing resonant flavour destroys it.
A second is generic: at $K=17$ the representation may compress \emph{any} axis
interior to a group, in which case the excess measures ``unlabelled versus
labelled'' and nothing about flavour. Both flavour axes lie interior to an
R16\_Q1 group and the topology axis does not, so these columns cannot separate
the accounts; the control in Section~4 can. Either way the observation stands:
L162 decodes QCD flavour with no QCD flavour label anywhere in its vocabulary, so
what a pretraining vocabulary erases is not confined to the jets whose labels it
merged.

\paragraph{The specific component is the better-determined one.} Across seeds
the difference-in-differences spans only \didRange{} against \rawRange{} for the
raw effect: what varies with the pretraining seed is the general degradation.
Taking the three coarse-arm seeds as the sampling unit (2 d.o.f.), the two
excesses are the seed-resolved quantities --- $t=\tDidRes$, $p=\pDidRes$ on
the merged axis and $t=\tDidQcd$, $p=\pDidQcd$ on the unlabelled one --- whereas
the general degradation of the retained axis is not ($t=\tTop$, $p=\pTop$). We
therefore report $\facTop\times$ as a point estimate whose sign is consistent
across all three seeds rather than as a resolved effect, and rest the paper's
claims on the excesses.

\paragraph{A nonlinear probe recovers part of it, not most of it.} On the
directly-labelled axis the MLP reduces R16\_Q1's $1-\mathrm{AUC}$ by \mlpRecovery{}
against \mlpRecoveryL{} for L162 --- and L162's gain is
$\Delta\mathrm{AUC}=9\times10^{-5}$, within its own AUC noise and of the opposite
sign on the topology axis --- so some of the lost information is present but
nonlinearly encoded. It does not close the gap: the difference-in-differences falls from
\didMean{} to \didMlpMean{}, still $\didMlpFactor\times$. Since a linear probe
only lower-bounds the mutual information, that residual is what distinguishes
degradation from a mere change of coordinates.

\paragraph{One seed on the fine arm.} No interval here carries a
pretraining-seed component. The three R16\_Q1 runs supply a scale: the largest
seed-to-seed contrast \emph{within} R16\_Q1 is $\seedAB$ \seedABCI, so the
vocabulary effect is $\seedFactorLo$ to $\seedFactorHi$ times it, and every
R16\_Q1 seed is worse than the single L162 seed on all three axes. That bounds
the concern rather than removing it, since the bound is measured on the coarse
arm and assumed to transfer.

\paragraph{A note on loss.} Cross-entropy over $K$ classes is a statement about
one partition of the data, so our arms' losses (\LossLsixtwo{} for L162,
\LossRsixteen{} for R16\_Q1) are not the same quantity, and their ordering
follows by construction from R16\_Q1's label being a function of L162's.
Normalising does not repair it: against the stream's label marginal
(\HyRsixteen{} nats for R16\_Q1, \HyLsixtwoLo{}--\HyLsixtwoHi{} for L162)
R16\_Q1 still has the lower loss, whereas against information captured its
\InfoRsixteen{} nats sit \emph{inside} L162's
\InfoLsixtwoLo{}--\InfoLsixtwoHi{} interval. Which arm learned more is not
decidable from one loss per arm. Scaling laws mapping pretraining loss to
downstream performance \citep{scaling} are fitted at a fixed vocabulary; the
hazard is on the reading side, when losses from different label spaces are
placed on one axis.

\paragraph{Pipeline validation.} On JetClass-I under an identical recipe our
pretrained arm reaches macro AUC \ePartOurs{} against the published \ePartPub{}
\citep{part}, so the comparison is not an artefact of the training path.

\section{What this does not show}

\paragraph{Only one fine-arm seed.} Section~3.1 bounds this against the coarse
arm's seed spread, but that bound is measured there and assumed to transfer. The
pretraining-seed SD of the probe contrast across the three coarse-arm seeds is
\seedSDres{}, \seedSDqcd{} and \seedSDtop{} in $\log(1-\mathrm{AUC})$ on the
three axes --- that, not a variance imported from another arm, is the scale.

\paragraph{How many, or which?} The central limitation. Our coarse arm reduces
$K$ from 162 to 17 \emph{and} erases specific physical distinctions, and this
design cannot separate the two. The discriminating experiment is a
semantics-matched random control: an arm at matched $K=17$ whose partition of the
\Nnative{} labels is random rather than physical. If it degrades the probed axis
as much, only capacity matters; if the physical contraction is worse,
\emph{which} distinctions a vocabulary keeps matters beyond \emph{how many}. We
pre-register it, and the prediction that the physical contraction will be worse.
\citet{chen} randomise the target while holding auxiliary labels fixed --- the
inverse manipulation.

\paragraph{Two rungs, not a curve.} \citet{mehta} report coarse feedback
sufficing for human-aligned visual representations, saturating near eight
categories; our coarse endpoint $K=17$ sits above that knee. The results do not
conflict --- they ask whether coarse supervision \emph{suffices} for a target the
coarse labels span, we ask what it \emph{destroys} elsewhere --- but a full curve
is needed to locate the boundary.

\paragraph{One transfer axis.} The QCD result rests on a single pair of native
labels. We have not shown the effect generalises to other unlabelled
distinctions, nor whether it is symmetric --- whether a vocabulary that merges
topology damages flavour the same way.

\section{Conclusion}

Varying only the label vocabulary, we find two effects: a general degradation
from coarsening, and on top of it the loss of an erased feature --- seen both
where the fine vocabulary labels it and where neither does. The label space is a
third axis of pretraining design, still set by convention, whose effects do not
stay where the labels changed.

\bibliographystyle{plainnat}
\bibliography{refs}

\end{document}
